\documentclass[12pt,letterpaper]{report}

% =========================================================
% PAQUETES BÁSICOS
% =========================================================
\usepackage[utf8]{inputenc}
\usepackage[T1]{fontenc}
\usepackage[spanish,es-tabla]{babel}
\usepackage{csquotes}
\usepackage{geometry}
\usepackage{graphicx}
\usepackage{float}
\usepackage{booktabs}
\usepackage{longtable}
\usepackage{array}
\usepackage{tabularx}
\usepackage{hyperref}
\usepackage{xcolor}
\usepackage{enumitem}
\usepackage{amsmath, amssymb}
\usepackage{caption}
\usepackage{subcaption}
\usepackage{url}
\usepackage{ifthen}

% =========================================================
% CONFIGURACIÓN DE PÁGINA
% =========================================================
\geometry{
    left=3cm,
    right=2.5cm,
    top=3cm,
    bottom=3cm
}

\hypersetup{
    colorlinks=true,
    linkcolor=blue,
    urlcolor=blue,
    citecolor=blue
}

\setlength{\parskip}{0.6em}
\setlength{\parindent}{0pt}

\renewcommand{\arraystretch}{1.2}

% =========================================================
% COMANDOS GENERALES
% =========================================================

% Marca texto pendiente. Antes de la entrega final no debe quedar ningún \pendiente{}.
\newcommand{\pendiente}[1]{\textcolor{red}{\textbf{[PENDIENTE: #1]}}}

% Entrada segura de capítulos. Si el archivo no existe, el documento compila con advertencia.
\newcommand{\inputchapter}[1]{
    \IfFileExists{#1.tex}{
        \input{#1}
    }{
        \chapter*{Archivo no encontrado}
        \addcontentsline{toc}{chapter}{Archivo no encontrado: #1}
        \textcolor{red}{No se encontró el archivo \texttt{#1.tex}.}
    }
}

% Tabla estándar para repositorios de software.
% Uso:
% \softwareRepo{Nombre}{URL}{Rama/versión/commit}{Responsable}{Función}{Entradas}{Salidas}{Estado}
\newcommand{\softwareRepo}[8]{
\begin{table}[H]
\centering
\caption{Repositorio asociado: #1}
\begin{tabularx}{\textwidth}{>{\bfseries}p{4cm} X}
\toprule
Nombre del software & #1 \\
Repositorio & \url{#2} \\
Rama, versión o commit & #3 \\
Responsable & #4 \\
Función dentro del sistema & #5 \\
Entradas & #6 \\
Salidas & #7 \\
Estado actual & #8 \\
\bottomrule
\end{tabularx}
\end{table}
}

% Tabla estándar de problemas, decisiones y ajustes.
% Los autores pueden copiar esta estructura y completar filas.
\newcommand{\problemTableTemplate}{
\begin{longtable}{p{2.6cm} p{2.6cm} p{2.6cm} p{3cm} p{3cm}}
\caption{Problemas presentados, decisiones y ajustes realizados.}\\
\toprule
\textbf{Problema} & \textbf{Evidencia} & \textbf{Causa probable} & \textbf{Decisión o ajuste} & \textbf{Resultado} \\
\midrule
\endfirsthead

\toprule
\textbf{Problema} & \textbf{Evidencia} & \textbf{Causa probable} & \textbf{Decisión o ajuste} & \textbf{Resultado} \\
\midrule
\endhead

\pendiente{Problema identificado} &
\pendiente{Log, captura, medición o prueba} &
\pendiente{Causa técnica probable} &
\pendiente{Cambio aplicado y justificación} &
\pendiente{Resultado posterior} \\

\bottomrule
\end{longtable}
}

% Tabla estándar de pruebas.
\newcommand{\testTableTemplate}{
\begin{longtable}{p{1.8cm} p{3.5cm} p{3cm} p{4cm} p{2.5cm}}
\caption{Matriz de pruebas del capítulo.}\\
\toprule
\textbf{ID} & \textbf{Prueba} & \textbf{Módulo} & \textbf{Resultado esperado} & \textbf{Estado} \\
\midrule
\endfirsthead

\toprule
\textbf{ID} & \textbf{Prueba} & \textbf{Módulo} & \textbf{Resultado esperado} & \textbf{Estado} \\
\midrule
\endhead

\pendiente{PR-XXX} &
\pendiente{Nombre de la prueba} &
\pendiente{Módulo} &
\pendiente{Resultado esperado} &
\pendiente{Pendiente/Aprobada/Fallida} \\

\bottomrule
\end{longtable}
}

% Tabla estándar de transferencia de conocimiento.
\newcommand{\knowledgeTransferTableTemplate}{
\begin{longtable}{p{4cm} p{5cm} p{5cm}}
\caption{Aprendizajes y conocimiento transferible.}\\
\toprule
\textbf{Tema} & \textbf{Aprendizaje obtenido} & \textbf{Recurso o evidencia asociada} \\
\midrule
\endfirsthead

\toprule
\textbf{Tema} & \textbf{Aprendizaje obtenido} & \textbf{Recurso o evidencia asociada} \\
\midrule
\endhead

\pendiente{Tema} &
\pendiente{Aprendizaje práctico} &
\pendiente{Figura, prueba, repositorio, tabla o evidencia} \\

\bottomrule
\end{longtable}
}

% =========================================================
% DATOS DEL DOCUMENTO
% =========================================================
\title{
    \textbf{Informe Final del Sistema}\\
    \large Documentación técnica, decisiones de diseño, pruebas y transferencia de conocimiento
}

\author{
    Equipo de desarrollo\\
    Martín Ramírez Espinosa\\
    Mateo Almeida Gómez\\
    David Ramírez Betancourth\\
    Oscar Andres Gutierrez Estepa\\ 
}

\date{\today}

% =========================================================
% INICIO DEL DOCUMENTO
% =========================================================
\begin{document}

\maketitle

\tableofcontents
\listoffigures
\listoftables

% =========================================================
% NOTA GENERAL PARA AUTORES
% =========================================================
\chapter*{Nota para los autores}
\addcontentsline{toc}{chapter}{Nota para los autores}

Este documento corresponde al \textbf{informe final del sistema}. El archivo \texttt{main.tex} es el archivo maestro y debe ser modificado únicamente por el integrador del documento.

Cada autor debe trabajar únicamente en el archivo correspondiente a su capítulo dentro de la carpeta \texttt{section/}. La estructura, reglas de escritura y criterios de calidad se encuentran en el archivo \texttt{README.md}.

Cada capítulo técnico debe documentar:

\begin{itemize}
    \item Propósito del capítulo.
    \item Contexto dentro del sistema general.
    \item Desarrollo técnico.
    \item Entradas, salidas e interfaces.
    \item Decisiones de diseño o implementación.
    \item Problemas presentados y ajustes realizados.
    \item Validación, pruebas o evidencias.
    \item Aprendizajes y transferencia de conocimiento.
    \item Limitaciones.
    \item Recomendaciones específicas.
\end{itemize}

Si se menciona software, el código fuente completo no debe incluirse en este informe. Debe explicarse la función del software y enlazarse el repositorio real correspondiente. Si el repositorio no existe todavía, debe escribirse explícitamente: \textit{Repositorio pendiente de publicación}.

Antes de la versión final no debe quedar ninguna marca \verb|\pendiente{}| en el documento.

% =========================================================
% CAPÍTULOS PRINCIPALES
% =========================================================

\inputchapter{section/1_introduction}
\inputchapter{section/2_requirements-restrictions}
\inputchapter{section/3_general_design_system}
\inputchapter{section/4_HW}
\inputchapter{section/5_SDR-acquisition}
\inputchapter{section/6_Platform_cloud}
\inputchapter{section/7_Spectral-localization_cloud}
\inputchapter{section/8_Test_protocols_edge-cloud}
\inputchapter{section/9_System_integration}
\inputchapter{section/10_General_conclusions}
\inputchapter{section/11_Recommendations}
\inputchapter{section/12_References}

% =========================================================
% ANEXOS
% =========================================================
\appendix

\inputchapter{section/A_IA-collaboration-technique}
\inputchapter{section/B_Additional_diagrams}
\inputchapter{section/C_Relevant_code_fragments}
\inputchapter{section/D_Manuals_guides}
\inputchapter{section/E_Test_results}

\end{document}
